Remote sensing has reshaped landslide susceptibility mapping by delivering globally consistent inputs: Sentinel-1 C-band SAR products for surface-deformation and amplitude features \citep{poggi2026sentinellandslide}, Sentinel-2 multi-spectral composites for surface cover and vegetation indices, MERIT and Copernicus 30~m DEMs for topographic and hydrologic derivatives, and WorldClim 2.1 bioclimatic surfaces for environmental context. These data, harmonized to a common grid, feed the dominant paradigm of fitting a machine-learning classifier to a per-basin inventory of labeled events---a workflow that produces strong results when local training data are abundant \citep{reichenbach2018review, merghadi2020ml, ma2025metalsm}. However, the practitioner facing a \emph{newly-studied} basin confronts a data-scarcity bottleneck: digitizing a usable inventory typically requires expert geomorphological mapping over hundreds of events, an effort that can take months and is rarely funded for individual watersheds. As a result, basins with sparse inventories ($N<50$ mass-movement records) cannot benefit from data-hungry deep learning approaches, while transfer of models trained on data-rich neighboring basins suffers from systematic shifts in geology, climate, and surface cover \citep{zhu2026thawsusceptibility}.

Two strategies have emerged to mitigate this. \textbf{Transfer learning} by fine-tuning a model pre-trained on data-rich basins \citep{ghorbanzadeh2022landslide} reduces label requirements but still needs tens to hundreds of target samples. \textbf{Domain adaptation} approaches such as DANN \citep{ganin2015dann} attempt to learn domain-invariant representations from labeled source data and unlabeled target data, with mixed success in remote sensing \citep{tuia2016domain, huang2026transferability}; recent work on source-free domain adaptation for cross-region Sentinel time-series mapping further documents how p(x) and p(y$\mid$x) shifts dominate cross-site transfer \citep{mohammadi2024sourcefreeunsupervised}.

A complementary paradigm, \textbf{meta-learning} \citep{finn2017maml, nichol2018reptile}, trains a model to \emph{adapt rapidly to new tasks from a few examples}. Meta-learning has shown promise in remote sensing for land-cover classification \citep{russwurm2020metalearning}, satellite image segmentation \citep{tseng2022timm}, and cross-scenario hyperspectral-to-Sentinel-2 trait estimation \citep{fu2025crossscenario}, and a recent single-site application \citep{ma2025metalsm} demonstrated temporal meta-learning for landslide susceptibility in Hong Kong. However, to our knowledge, no published study has systematically evaluated meta-learning for \emph{cross-basin} landslide susceptibility transfer, particularly across the extreme hyper-arid-to-hyper-humid climate gradient of Chile, which spans nearly 30$^\circ$ of latitude and where traditional transfer learning is difficult. Spatially explicit cross-validation protocols, essential for unbiased ML evaluation on autocorrelated Earth-observation data \citep{brenning2005spatial}, remain under-applied in landslide remote sensing.

In this work, we present the first systematic \emph{cross-basin} meta-learning benchmark for landslide susceptibility classification across a hydro-climatic gradient in north-central Chile, complementing the temporal single-site setting of \citet{ma2025metalsm}. Our contributions are:

\begin{enumerate}
    \item We construct a unified, balanced dataset spanning \textbf{14 Chilean basins} (10 source, 4 target) with consistent feature engineering: 30~m DEM derivatives (terrain, hydrology, texture), Sentinel-2 spectral bands, WorldClim bioclimatic variables, and SERNAGEOMIN geology categoricals. Inventory labels span four published catalogs and one ad-hoc digitization (3{,}290 positives + 3{,}290 negatives after balanced sampling).
    \item We benchmark \textbf{five methods}---Independent (no source), Fine-tune (concat-source pretrain), Reptile, FOMAML, DANN---plus \textbf{four classical machine-learning baselines} (logistic regression, random forest, XGBoost, CatBoost), across $K$-shot evaluation ($K \in \{1, 5, 10, 20\}$) on four target basins, with three random seeds and bootstrap 95\% confidence intervals (62{,}520 model fits in total).
    \item We report results under \emph{both} random and spatial cross-validation protocols, addressing the spatial autocorrelation concern that pervades landslide ML evaluation \citep{brenning2005spatial}.
    \item Our findings: (i) meta-learning (Reptile or FOMAML) wins $K=1$ over Independent in 4/4 target basins under spatial CV, with a mean lift of $+5.5$\,pp $F_1$ (range $+4.3$ to $+7.8$\,pp); Reptile and FOMAML perform within overlapping confidence intervals in the majority of target-$K$ cells; (ii) the advantage decays monotonically with $K$ and vanishes by $K=20$; (iii) DANN does not improve over Independent on average and degrades substantially in the smallest target basin (Copiap\'o $N=19$: $-6$\,pp at $K=1$), suggesting adversarial alignment is counter-productive when both source domains and target inventories are small; (iv) adaptation curves confirm meta-learned initializations reach $F_1 \approx 0.78$ within $\leq 3$ gradient steps in Copiap\'o versus $>100$ steps for Independent; (v) classical ML achieves competitive $F_1$ under spatial CV, suggesting meta-learning's principal benefit is rapid adaptation efficiency rather than peak performance.
\end{enumerate}

The remainder of this paper is organized as follows. Section~\ref{sec:related} reviews related work in landslide ML, transfer learning, and meta-learning. Section~\ref{sec:data} describes the study area and data. Section~\ref{sec:methods} details the methodology. Section~\ref{sec:results} presents results. Section~\ref{sec:discussion} discusses implications and limitations. Section~\ref{sec:conclusions} concludes.
